Considera el punt $P(3,-2,4)$ i el pla $\pi\colon x-2y+2z-6=0$.

\begin{apartats}

\apartat[1,25]{0,75}
Troba la projecció ortogonal $Q$ del punt $P$ sobre el pla $\pi$.

\begin{solucio}
La recta perpendicular a $\pi$ per $P$ és $(3+t,\ -2-2t,\ 4+2t)$. Substituint al pla:
$(3+t)-2(-2-2t)+2(4+2t)-6=9+9t=0$, $t=-1$. La projecció és $Q=(2,0,2)$.
\end{solucio}

\begin{tria}{simetria}
\itemtria{simetric-pla}{1}{1,25}
Troba el punt simètric $P'$ de $P$ respecte del pla $\pi$.

\begin{solucio}
$Q$ és el punt mitjà de $P$ i $P'$: $P'=2Q-P=(4-3,\ 0+2,\ 4-4)=(1,2,0)$.\\
Comprovació: $\overrightarrow{PP'}=(-2,4,-4)$ és proporcional a $\vec n=(1,-2,2)$, i el punt mitjà,
$(2,0,2)$, és de $\pi$.
\end{solucio}

\itemtria{pla-mediador}{1}{1,25}
Troba el pla respecte del qual els punts $A(1,2,3)$ i $B(3,6,-1)$ són simètrics.

\begin{solucio}
És el pla perpendicular a $\overrightarrow{AB}=(2,4,-4)$, o a $(1,2,-2)$, que passa pel punt mitjà de
$A$ i $B$, $M(2,4,1)$: $x+2y-2z+D=0$ amb $2+8-2+D=0$, $D=-8$. El pla és $x+2y-2z-8=0$.
\end{solucio}
\end{tria}

\begin{nomesllarg}
\apartat{0,75}
Calcula la distància de $P$ a $\pi$ amb la fórmula, i comprova que és igual a $|\overrightarrow{PQ}|$.

\begin{solucio}
$d(P,\pi)=\dfrac{|3+4+8-6|}{\sqrt{1+4+4}}=\dfrac93=3$, i $\overrightarrow{PQ}=(-1,2,-2)$ té mòdul
$\sqrt{1+4+4}=3$.
\end{solucio}
\end{nomesllarg}

\end{apartats}
